\section{El problema completo, antes de entrar en los cálculos}\label{sec:mapa}

\subsection{La decisión que queremos tomar}
Kaido debe decidir cuánto encargar a fábrica sin conocer las ventas futuras.
Si encarga poco, puede perder compradores; si encarga demasiado, mantiene coches
que cuestan dinero y envejecen. Además, el pedido no llega de inmediato. Una
decisión tomada después de observar una subida de ventas tarda tres meses en
convertirse en una llegada modificable a la red.

El objetivo del trabajo no es producir una cifra de ventas aislada. Es construir
una \textbf{regla de decisión que siga funcionando cuando la previsión se
equivoca}. Por eso se necesitan dos herramientas diferentes:
\begin{enumerate}
  \item Una \textbf{previsión}: estima cuántos clientes querrán comprar cada mes.
  \item Una \textbf{simulación}: representa qué puede venderse realmente con el
  stock disponible, los compromisos de suministro y la regla de revisión.
\end{enumerate}
La recomendación económica se obtiene al comparar políticas sobre las mismas
trayectorias de demanda. El modelo trabaja por mes y por segmento; no representa
individualmente los 64 concesionarios.

\begin{figure}[htbp]
\centering
\begin{tikzpicture}[node distance=5mm,
  caja/.style={draw=gris!40,fill=fondo,rounded corners=1pt,text width=3.2cm,
    minimum height=1.65cm,align=center,font=\small},
  flecha/.style={-{Stealth[length=2mm]},draw=azul,thick}]
  \node[caja] (d) {\textbf{Datos}\\Qué ocurrió\\Clientes y episodios};
  \node[caja,right=of d] (p) {\textbf{Previsión}\\Qué podría ocurrir\\Mercado y posición};
  \node[caja,right=of p] (s) {\textbf{Simulación}\\Qué podremos atender\\Stock y plazos};
  \node[caja,right=of s] (e) {\textbf{Decisión}\\Qué política elegir\\Servicio y margen};
  \draw[flecha] (d)--(p);\draw[flecha] (p)--(s);\draw[flecha] (s)--(e);
\end{tikzpicture}
\caption{Las cuatro piezas de la solución. Un error en una fase se transmite a
las siguientes: una demanda mal interpretada puede terminar en un pedido mal dimensionado.}
\label{fig:proceso}
\fuente{estructura de los cuatro notebooks revisados}
\end{figure}

\subsection{Cómo se organiza el proyecto}
\begin{table}[htbp]\centering
\caption{Responsabilidad y resultado de cada notebook.}\label{tab:mapa-notebooks}
\begin{tabularx}{\linewidth}{@{}p{3.4cm}YY@{}}\toprule
Notebook & Pregunta principal & Resultado que entrega \\\midrule
\archivo{datos.ipynb} & ¿Qué parte del histórico representa compras de clientes?
& Series observadas, diagnóstico y ajustes de episodios.\\
\archivo{demanda.ipynb} & ¿Qué demanda mensual podemos prever y qué error tiene?
& Previsión, evaluación temporal y parámetros de incertidumbre.\\
\archivo{simulacion.ipynb} & ¿Cómo se transforman pedidos y llegadas en ventas y costes?
& Simulador comprobado y referencia determinista.\\
\archivo{escenarios.ipynb} & ¿Qué política compensa y qué riesgos mantiene?
& Simulaciones, política elegida, sensibilidades y entrega.\\\bottomrule
\end{tabularx}
\end{table}

Los notebooks comparten las celdas de definiciones etiquetadas \archivo{api}.
Las fases posteriores cargan esas funciones en un espacio independiente; no
ejecutan de nuevo los gráficos de otra fase ni usan las tablas antiguas como si
fueran el modelo. \archivo{analisis.ipynb} y los módulos históricos de previsión
y simulación se conservan como referencia inicial. Este informe describe la
\textbf{versión revisada}.

\subsection{Tres niveles de afirmación}
Conviene identificar en cada resultado de dónde procede la cifra:
\begin{description}[style=nextline]
  \item[Observado.] Se lee en los datos o se calcula directamente a partir de
  ellos: por ejemplo, matriculaciones menos automatriculaciones.
  \item[Estimado o supuesto.] Requiere un método o una hipótesis: por ejemplo,
  cuántas compras se adelantaron o cuántos coches vienen ya en camino.
  \item[Simulado.] Es el resultado de aplicar reglas y supuestos a una posible
  trayectoria futura: ventas, stock y contribución de una réplica.
\end{description}
La distinción importa porque repetir mil veces una simulación no convierte sus
supuestos en hechos. Tampoco una correlación histórica demuestra que una variable
sea la causa de un cambio en las ventas.

\subsection{Vocabulario mínimo para leer el resto}
\begin{table}[htbp]\centering
\caption{Conceptos que no deben confundirse.}\label{tab:glosario-inicial}
\begin{tabularx}{\linewidth}{@{}p{3.2cm}Y@{}}\toprule
Concepto & Interpretación en este trabajo \\\midrule
Matriculaciones & Incluyen ventas a clientes y automatriculaciones.\\
Venta a cliente & Compra observada después de descontar automatriculaciones.\\
Demanda potencial & Lo que los clientes querrían comprar; puede superar las ventas si falta stock.\\
Wholesale / retail & Facturación a la red / venta al cliente, respectivamente.\\
PFF & Precio franco fábrica; no es el PVP con impuestos.\\
VM & Gasto comercial: descuentos, campañas e incentivos por vehículo.\\
CVP & Contribución del vehículo después del VM: beneficio bruto menos VM.\\
Transfer & Vehículo procedente del stock de otro país europeo para cubrir un faltante.\\
FIFO & Se vende primero la cohorte más antigua disponible.\\
Escenario & Conjunto coherente de supuestos sobre A013.\\
Réplica & Una trayectoria aleatoria completa dentro de un escenario.\\\bottomrule
\end{tabularx}
\fuente{enunciado, secciones 2 y 9, y convenciones de la simulación revisada}
\end{table}

\begin{idea}{La idea que conviene conservar}
Una previsión de \num{1000} coches no significa que se venderán \num{1000}.
Primero tiene que existir esa demanda y después tiene que haber coches para
atenderla. Por último, vender esos coches debe evaluarse con sus costes: más
unidades no siempre implican más contribución.
\end{idea}
\FloatBarrier
